\chapter*{摘要}
\addcontentsline{toc}{chapter}{摘要}
\markboth{摘要}{摘要}

本作品在正点原子领航者 ZYNQ7020 开发板上实现集成 Ascon-AEAD128 的 RISC-V SoC。 PL 内为五级流水线 RV32IM + Zicsr 处理器，包含分支预测、I/D Cache、AXI 互连及 UART、GPIO、CLINT。密码功能以 NIST SP 800-232 Final 为基准，提供纯软件、存储器映射加速器和定制指令三种方案。四条定制指令使用 custom-0 编码空间，以 inline assembly 调用，不修改 GCC backend。

当前 PYNQ 发布版由 PS 经两个 AXI-Lite 窗口及 PL 内 UART mailbox 控制 RV32 服务固件。明文和关联数据各不超过 16 B，三模式各通过 289 组官方容量内 KAT 及 289 次错误 Tag 拒绝。 Notebook 顺序执行、9 项错误恢复门禁与每模式 50 次加解密往返均已有实板记录。固定官方 KAT \#529（PT=16 B、AD=0），各预热 5 次后采样 30 次，加密 SW/HW/ISE 中位周期为 23342/1439/3995，HW 与 ISE 的周期加速比为 16.221$\times$、5.843$\times$；解密为 24005/1388/4385 周期。周期来自固件 AEAD 调用边界，不含 PS 与串口；PS 调用墙钟中位 4.262--5.216 ms 单列，不充作同倍数响应加速。

Vivado 2025.2 对当前 \texttt{pynq\_\allowbreak{}wrapper} 发布版在配置 100 MHz 下实现 WNS +0.034 ns、TNS 0、 setup/hold 失败端点 0、DRC 0 Error/0 Critical Warning，普通 Warning 如实保留。另实现独立的 256 路 RO-PUF 与 Device DNA 密钥通路原型，同次上电的注册/重建自检通过，最新有效日志注册质量为 128/128；不声称跨上电稳定、密钥保密隔离或完整安全启动。 host 完整 1089 组 KAT 与板级 289 组分列；独立频率、小时级稳定性与干净环境端到端复现尚未完成。

\textbf{关键词}：RISC-V；Ascon-AEAD128；定制指令扩展；MMIO；Zynq-7000；PYNQ；RO-PUF 原型
